%==============================================================================%
% Chapter 2 – Background
%==============================================================================%

\chapter{Background}
\label{ch:background}

\section{Lattice-Based Cryptography}

\subsection{Module Lattices}

A \emph{lattice} $\Lambda \subseteq \Z^n$ is a discrete additive subgroup of
$\Z^n$; equivalently, it is the $\Z$-span of a basis matrix $\mathbf{B} \in
\Z^{n \times n}$.  The \emph{shortest vector problem} (SVP) asks for the
shortest non-zero lattice vector under the Euclidean norm, and is conjectured
to be hard for quantum computers.

Module lattices generalise to modules over a polynomial ring.  Let $R = \Z[X]
/ (X^n + 1)$ and $R_q = R/qR$.  A \emph{module lattice} of rank $k$ over $R$
is a free $R$-module $M \cong R^k$ together with a $\Z$-embedding into $\Z^{kn}$.
The \emph{Module Short Integer Solution} (MSIS) and \emph{Module Learning With
Errors} (MLWE) problems over such lattices underpin the security of Dilithium,
FALCON, and \haetae{}.

\subsection{NTT-Friendly Parameters}

For efficient computation, the ring $R_q$ must support a fast polynomial
multiplication algorithm.  If $q \equiv 1 \pmod{2n}$ then $X^n + 1$ splits
completely over $\F_q$ and the NTT reduces polynomial multiplication to
componentwise multiplication of length-$n$ vectors.  If only $q \equiv 1
\pmod{n}$ (but $q \equiv 3 \pmod{4}$ or similar), a \emph{negacyclic} NTT can
still be defined~\cite{Longa2016}.

For \haetae{}, $n = 256$ and $q = 64513 \equiv 1 \pmod{512}$.  This means a
primitive $512$th root of unity $\zeta_0$ exists in $\F_q$, enabling a full
$256$-point negacyclic NTT.

\section{Number-Theoretic Transform}

\subsection{Definition}

Let $q$ be an odd prime with $q \equiv 1 \pmod{2n}$, and let $\zeta \in \F_q^*$
be a primitive $2n$th root of unity (so $\zeta^{2n} = 1$, $\zeta^n = -1$).
The \emph{negacyclic NTT} of a polynomial $f = \sum_{j=0}^{n-1} f_j X^j \in
R_q$ is the vector $\hat{f} \in \F_q^n$ defined by
\begin{equation}
  \hat{f}_i = \sum_{j=0}^{n-1} f_j \cdot \zeta^{(2\brev_8(i)+1) \cdot j},
  \qquad i = 0, \dots, n-1,
\label{eq:ntt_def}
\end{equation}
where $\brev_8 : \{0,\dots,255\} \to \{0,\dots,255\}$ is the bit-reversal
permutation of an 8-bit index.

The inverse NTT is
\begin{equation}
  f_i = \frac{1}{n} \sum_{j=0}^{n-1} \hat{f}_j \cdot \zeta^{-((2\brev_8(j)+1) \cdot i)},
  \qquad i = 0, \dots, n-1.
\label{eq:intt_def}
\end{equation}

\begin{proposition}[Round-trip identity]
  $\INTT(\NTT(f)) = f$ for all $f \in R_q$.
\label{prop:roundtrip}
\end{proposition}

\subsection{Cooley--Tukey Butterfly Algorithm}
\label{sec:ct_butterfly}

The NTT defined by \cref{eq:ntt_def} can be computed in $O(n \log n)$ field
operations using the Cooley--Tukey decimation-in-time (DIT) algorithm.  For
the negacyclic case with $n = 256 = 2^8$, the algorithm proceeds in 8 stages.

At stage $s$ (with $s = 0, 1, \dots, 7$, from outermost to innermost), the
half-length is $\ell = 2^{7-s}$.  If $g$ is the group number within that
stage, then the implementation's table counter is $k = 2^s + g$ and the
twiddle factor is
\[
  w = \zeta^{\brev_8(k)}.
\]
Index 0 of the concrete forward table is padding, so the loop consumes
indices $1,\dots,255$.  The butterfly operation on a pair $(a, b)$ with
twiddle factor $w$ is:
\begin{equation}
  (a', b') = (a + w \cdot b,\ a - w \cdot b).
\label{eq:butterfly}
\end{equation}

\begin{algorithm}[htbp]
\caption{Forward NTT (Cooley--Tukey DIT, negacyclic, $n = 256$)}
\label{alg:ntt}
\begin{algorithmic}[1]
\Procedure{NTT}{$a[0..255]$}
\State $k \gets 1$
\For{$\ell = 128, 64, 32, 16, 8, 4, 2, 1$}
  \For{$\mathit{start} = 0, 2\ell, 4\ell, \dots, 254$}
    \State $w \gets \zeta^{\brev_8(k)}$ \Comment{$k$-th twiddle factor}
    \State $k \gets k + 1$
    \For{$j = \mathit{start}, \dots, \mathit{start}+\ell-1$}
      \State $t \gets w \cdot a[j+\ell] \bmod q$
      \State $a[j+\ell] \gets a[j] - t$
      \State $a[j] \gets a[j] + t$
    \EndFor
  \EndFor
\EndFor
\EndProcedure
\end{algorithmic}
\end{algorithm}

\subsection{Negacyclic NTT vs.\ Standard NTT}

A standard (cyclic) DFT of length $n$ uses the recurrence for the cyclic ring
$\Z_q[X]/(X^n - 1)$.  The \emph{negacyclic} ring $\Z_q[X]/(X^n + 1)$ requires
a pre-twist by powers of $\zeta$ (a half-integer shift of the exponent), which
accounts for the $+1$ in the exponent formula of \cref{eq:ntt_def}.

Critically for this dissertation, \haetae{} uses a \emph{complete} $n=256$
negacyclic NTT with bit-reversal exponent $\brev_8(k)$, whereas MLKEM uses an
incomplete $128$-point transform with exponent $\brev_7(k)$.  This distinction
has direct consequences for which algebraic lemmas can be shared between
verification projects.

\section{Montgomery Modular Arithmetic}

\subsection{REDC Algorithm}

Montgomery reduction~\cite{Montgomery1985} is a technique for efficient modular
multiplication that avoids expensive division.  Given $a \in \Z$ with
$0 \le a < q \cdot R$ and integers $q, R$ with $\gcd(q, R) = 1$, the
Montgomery product is $a \cdot R^{-1} \bmod q$.

\begin{algorithm}[htbp]
\caption{Montgomery Reduction with the \haetae{} Sign Convention}
\label{alg:mont_reduce}
\begin{algorithmic}[1]
\Require $a, R = 2^{32}$, $q^+ = q^{-1} \bmod R$
\Ensure $a \cdot R^{-1} \bmod q$
\State $m \gets (a \bmod R) \cdot q^+ \bmod R$
\State $t \gets (a - m \cdot q) \div R$
\If{$t < 0$} \Return $t + q$ \Else\ \Return $t$ \EndIf
\end{algorithmic}
\end{algorithm}

In the \haetae{} implementation, $R = 2^{32}$, $q = 64513$,
$\mathsf{QINV} = 940508161$ ($\equiv q^{-1} \bmod 2^{32}$), and
$\mathsf{MONT} = R \bmod q = 14321$.  Thus the implementation uses the
positive-inverse convention in \Cref{alg:mont_reduce}: it computes
$m = (a \bmod R)\mathsf{QINV} \bmod R$ and then forms $a - m q$.

\subsection{Montgomery Multiplication}

Given $a, b \in \Z_q$ with both represented in Montgomery form (i.e., as
$a \cdot R \bmod q$ and $b \cdot R \bmod q$), their product in Montgomery form
is computed as $\mathrm{REDC}(a \cdot b)$, since
$\mathrm{REDC}(aR \cdot bR) = ab R^2 \cdot R^{-1} = abR \bmod q$.

\subsection{Signed Montgomery Reduction}

The \haetae{} implementation uses a \emph{signed} variant where intermediate
values are 32-bit signed integers and the final conditional subtraction is
replaced by an arithmetic right shift.  Specifically, for $a \in \Z$ with
$\abs{a} < q \cdot 2^{15}$:
\begin{equation}
  a \;-\; \text{int32}(a_{\text{lo}} \cdot \mathsf{QINV}) \cdot q
  \quad\text{(as signed 64-bit)}
  \;\text{ shifted right by 32}
\label{eq:signed_mont}
\end{equation}
yields a value congruent to $a \cdot R^{-1} \pmod{q}$ and, in the checked
signed REDC theorem, a representative in $[-q,q)$.

\section{The EasyCrypt Proof Assistant}

\subsection{Architecture}

\easycrypt{}~\cite{Barthe2011} is an interactive proof assistant specialised
for reasoning about cryptographic programs.  It combines:
\begin{itemize}
  \item A \emph{program logic} layer supporting Hoare logic (HL) for
        deterministic programs and probabilistic relational Hoare logic
        (pRHL) for pairs of probabilistic programs.

  \item An \emph{ambient logic} layer that is a form of higher-order logic
        (HOL) with a type theory supporting algebraic structures (groups,
        rings, fields).

  \item A tactic engine with automation backends including SMT solvers
        (CVC4, Z3) and a computer algebra system for ring/field equalities.
\end{itemize}

\subsection{Hoare Logic in EasyCrypt}

A \emph{Hoare triple} $\{P\}\; c\; \{Q\}$ asserts that if precondition $P$
holds before executing command $c$, then postcondition $Q$ holds after.  In
\easycrypt{}, Hoare triples over \emph{procedures} (modules with state) are
written \ec{hoare[M.f : P ==> Q]}.

A \emph{relational Hoare triple} $\{P\}\; (c_1 \sim c_2)\; \{Q\}$ asserts
that if $c_1$ and $c_2$ start in states satisfying $P$, their outputs satisfy
$Q$.  This is the foundation for proving that two programs are equivalent.

\subsection{Module System}

\easycrypt{} organises code into \emph{modules}.  A module interface (type) is
a signature listing procedure names and types; a module is a concrete
implementation.  The \jasmin{} compiler generates an \easycrypt{} module for
each exported function, which can then be related to hand-written specification
modules via pRHL.

\subsection{Theory Cloning}

Algebraic structures are developed as parameterised \emph{theories} that are
instantiated by \emph{cloning}: \ec{clone Ring.ComRing as Zq with type t = int, ...}
creates a copy of the commutative ring theory with $\Z_q$ as the carrier type
and appropriate witnesses for the ring axioms.

\section{The Jasmin Programming Language}

\jasmin{}~\cite{Almeida2017} is a C-like language designed for high-assurance
cryptographic code.  Its distinguishing features are:

\begin{itemize}
  \item \textbf{Explicit register allocation.}  Variables are annotated
        \jasm{reg} (in a register), \jasm{stack} (on the stack), or
        \jasm{ptr} (a pointer).  The compiler enforces the annotation.

  \item \textbf{Data types matching machine types.}  The types
        \jasm{u32}, \jasm{u64}, \jasm{u256} correspond directly to
        32-bit, 64-bit, and 256-bit (AVX2) machine words.

  \item \textbf{Safety verification.}  A safety checker ensures no
        undefined behaviour (array out of bounds, use before initialisation,
        integer overflow in unsigned arithmetic).

  \item \textbf{Constant-time verification.}  The companion
        \texttt{jasmin-ct} checker verifies that secret values do not control
        branches or memory addresses.  Its speculative mode also tracks
        transient pointer inputs, using explicit \jasm{\#[ct = "..."]} and
        \jasm{\#[sct = "..."]} contracts on exported procedures.

  \item \textbf{EasyCrypt extraction.}  Each \jasm{export fn} is compiled
        to both assembly and an \easycrypt{} module.  The extracted module has
        a semantic identical to the assembly, by the Jasmin compiler
        correctness theorem.
\end{itemize}

\begin{lstlisting}[language=Jasmin, caption={Jasmin function signature example}]
export fn poly_ntt_jazz(reg mut ptr u32[256] rp)
  -> reg ptr u32[256]
{
  rp = _poly_ntt(rp);
  return rp;
}
\end{lstlisting}

The export annotation triggers both assembly generation and \easycrypt{}
extraction.  The extracted \easycrypt{} type for \jasm{ptr u32[256]} is a
\ec{BArray1024.t} byte-array value, with coefficients accessed by
\ec{BArray1024.get32} and \ec{BArray1024.set32}.

\section{Related Work}

\paragraph{MLKEM/Kyber NTT.}
The closest predecessor is the Jasmin/EasyCrypt verification of the NTT in
ML-KEM~\cite{Barbosa2023}.  That work verifies the incomplete $128$-point
Cooley--Tukey NTT over a related ring structure but with an incomplete
transform and a different twiddle-indexing discipline.  Our work adapts
techniques from that development but requires new algebra lemmas for the
full $256$-point transform.

\paragraph{Dilithium NTT.}
The EasyCrypt verification of Dilithium~\cite{Barbosa2021} also formalises
an NTT in $\Z_q[X]/(X^{256}+1)$ but uses a different prime $q = 8380417$ with
different twiddle factor structure.

\paragraph{FALCON.}
FALCON uses a floating-point NTT over $\Z[X]/(X^n+1)$~\cite{Fouque2018}.
Its Jasmin implementation has been verified using \easycrypt{} but the
field arithmetic is quite different from the integer case studied here.

\paragraph{High-assurance cryptography.}
The HACL* project~\cite{Zinzindohoue2017} produces C code verified in
F*/Low*, with machine-checked compilation to assembly.  While HACL* has
verified ChaCha20 and Poly1305, it has not (at time of writing) produced a
verified lattice-based NTT.  Our work contributes to the complementary
Jasmin/EasyCrypt ecosystem, which has better support for the bitwise
operations prevalent in lattice primitives.
